\PaperPage{7}{1.1 Hilbert triples and conventions}
\label{hilb:triples}
The operator language begins with a complex Hilbert space $\mathcal H$,
a bounded linear operator $T\colon\mathcal H\to\mathcal H$, and a
distinguished vector $v$. Throughout, the inner product is
conjugate-linear in its first argument and linear in its second:
\[
 \langle ax,by\rangle=\overline a\,b\langle x,y\rangle,
 \qquad \|x\|^2=\langle x,x\rangle.
\]
This convention fixes the signs of the resolvent identities below. The
zero vector gives an identically zero scalar function; normalization
will usually mean $\|v\|=1$, but the elementary estimates retain the
factor $\|v\|^2$ so that weighted vectors are included.

\begin{definition}
A \emph{contractive Hilbert-vector model} is a triple
$(\mathcal H,T,v)$ with $\|T\|\le1$. It is \emph{positive} when $T$ is
self-adjoint and $0\le T\le I$. Its scalar orbit coefficients are
\[
 c_n=\langle v,T^nv\rangle,\qquad n\ge0.
\]
\end{definition}
Positivity in this definition is operator positivity, not positivity
of individual matrix entries in a chosen basis. In particular, it is
invariant under unitary change of coordinates. The construction also
singles out a vector: the operator alone does not specify the scalar
coefficients.

The natural observable subspace is
\[
 \mathcal H_v=\overline{\operatorname{span}}
       \{v,Tv,T^2v,\ldots\}.
\]
It is invariant under $T$ by boundedness. If $T=T^*$, then it is reducing.
Indeed, if $x\perp\mathcal H_v$ and $y\in\mathcal H_v$, then
$\langle y,Tx\rangle=\langle Ty,x\rangle=0$, since
$Ty\in\mathcal H_v$. Thus $\mathcal H_v^\perp$ is invariant as well.
The vector is called \emph{cyclic} if $\mathcal H_v=\mathcal H$.

For self-adjoint $T$, every scalar moment is therefore already a
moment of the restriction to $\mathcal H_v$:
\[
 \langle v,T^nv\rangle
 =\langle v,(T|_{\mathcal H_v})^nv\rangle.
\]
The complement is invisible to these data. This observation will later
prevent an unwarranted conclusion about the entire spectrum of an
operator reconstructed from one scalar function. These Hilbert-space
definitions are classical mathematical background; they are stronger
than the minimal Gram-orbit interface used in the audited source.

\PaperPage{8}{1.2 Coefficient bounds and the two defects}
\label{hilb:defects}
Let $(\mathcal H,T,v)$ be a contractive model. Iteration preserves the
norm bound, so $\|T^nv\|\le\|v\|$ for every $n$. Combining this fact
with Cauchy--Schwarz immediately bounds the scalar orbit:
\[
 |c_n|=|\langle v,T^nv\rangle|
       \le\|v\|\,\|T^nv\|\le\|v\|^2.
\]
It is useful to retain the two reasons for this bound separately.
Write $x=T^nv$, and define the contraction and Gram defects by
\[
 A=\|v\|^2(\|v\|^2-\|x\|^2),\qquad
 B=\|v\|^2\|x\|^2-|\langle v,x\rangle|^2.
\]
Contractivity gives $A\ge0$, whereas Cauchy--Schwarz gives $B\ge0$.
The following identity is purely algebraic.

\begin{proposition}\label{hilb:defect-identity}
For every $n\ge0$,
\[
 \|v\|^4-|c_n|^2=A+B.
\]
In particular, the coefficient defect is a sum of two nonnegative
quantities. Equivalently, it is the Euclidean norm-square of
$(\sqrt A,\sqrt B)\in\mathbb R^2$.
\end{proposition}
\begin{proof}
Add the displayed definitions of $A$ and $B$. The two occurrences of
$\|v\|^2\|x\|^2$ cancel. Nonnegativity justifies the square roots and
the norm-square description.
\end{proof}

This decomposition separates two possible losses. The orbit may lose
norm under iteration, producing a positive contraction defect. Even
when its norm is preserved, the orbit vector may cease to be aligned
with the distinguished vector, producing a positive Gram defect. No
self-adjointness is needed for either statement.

If $v\ne0$, equality $|c_n|=\|v\|^2$ occurs precisely when
$T^nv=\lambda v$ for some scalar $\lambda$ with $|\lambda|=1$.
To see this, equality in the coefficient bound forces both defects to
vanish. Equality in Cauchy--Schwarz makes $x$ a scalar multiple of
$v$, and equality of norms determines the modulus of that scalar.
The converse follows by substitution.

The source development isolates the real counterpart of this argument
using a norm-square function and a squared Cauchy inequality
\cite{source}. Such an interface is sufficient for the coefficient
estimate. It does not, by itself, construct a Hilbert space or prove
linearity of its orbit map.

\PaperPage{9}{1.3 Neumann series and disk analyticity}
\label{hilb:neumann}
For $|z|<1$, the operator series
\[
 R_z=\sum_{n=0}^{\infty}z^nT^n
\]
converges in operator norm, because $\|z^nT^n\|\le|z|^n$. Multiplying
its finite partial sums by $I-zT$ gives
\[
 (I-zT)\sum_{n=0}^{N}z^nT^n
 =I-z^{N+1}T^{N+1}.
\]
The remainder tends to zero in norm. The same calculation on the
other side proves $R_z=(I-zT)^{-1}$ and
$\|R_z\|\le(1-|z|)^{-1}$.

\begin{proposition}\label{hilb:disk-resolvent}
The scalar generating function of a contractive model is holomorphic
on the unit disk and satisfies
\[
 F(z)=\sum_{n=0}^{\infty}c_nz^n
     =\langle v,(I-zT)^{-1}v\rangle,\qquad
 |F(z)|\le\frac{\|v\|^2}{1-|z|}.
\]
\end{proposition}
\begin{proof}
The map $A\mapsto\langle v,Av\rangle$ is a continuous linear
functional on bounded operators, of norm at most $\|v\|^2$.
Apply it to the operator-norm convergent Neumann series.
On each closed disk $|z|\le r<1$, the geometric estimates are uniform,
so the series is holomorphic there and hence throughout the open disk.
\end{proof}

The argument also controls derivatives. Since $T$ commutes with
$I-zT$, it commutes with its inverse. Differentiating
$(I-zT)R_z=I$ gives $R_z'=R_zTR_z$. Induction consequently yields
\[
 F^{(k)}(z)=k!\,
   \langle v,T^k(I-zT)^{-k-1}v\rangle,
 \qquad
 |F^{(k)}(z)|\le
   \frac{k!\|v\|^2}{(1-|z|)^{k+1}}.
\]
In particular, $F^{(k)}(0)=k!c_k$. The coefficient sequence is
therefore the actual Taylor sequence of this analytic function, not
merely a formal expression assigned to the model.

The same majorant gives a quantitative approximation by finite data:
for $|z|\le r<1$,
\[
 \left|F(z)-\sum_{n=0}^Nc_nz^n\right|
 \le\frac{\|v\|^2r^{N+1}}{1-r}.
\]
This estimate is uniform in the model once the vector norm is fixed.
It controls truncation inside the disk, without supplying control
across its boundary.

The unit disk is a guaranteed domain, not necessarily the maximal
one. For example, a strict contraction gives a larger Neumann disk,
and spectral information can provide continuation far beyond it.
The construction at this stage uses neither a positive measure nor
an arithmetic calibration. Those additional structures will be needed
to obtain a half-plane sign and an identity with a prescribed function.

\PaperPage{10}{1.4 Positive contractions and complete differences}
\label{hilb:positive}
Assume now that $T=T^*$, $0\le T\le I$, and $\|v\|=1$.
The order inequalities mean
$0\le\langle x,Tx\rangle\le\|x\|^2$ for every $x$.
The spectrum of $T$ lies in $[0,1]$. The continuous functional
calculus for bounded self-adjoint operators preserves real
nonnegative functions on that interval \cite{conway}.

Put $u_n=\langle v,T^nv\rangle$. Positivity of $T^n$ and
$I-T^n$ gives $0\le u_n\le1$, and $u_0=1$. More generally,
the commuting operators $T$ and $I-T$ satisfy
\[
 \sum_{r=0}^{k}(-1)^r\binom kr u_{n+r}
 =\langle v,T^n(I-T)^kv\rangle\ge0.
\]
The final inequality follows by applying functional calculus to
$q^n(1-q)^k$, which is nonnegative for $0\le q\le1$.
These finite-difference inequalities contain more information than
the bound $|u_n|\le1$.

\begin{proposition}\label{hilb:bernstein-operator}
The Bernstein row associated with the moments of a positive model is
a nonnegative row of mass one:
\[
 w_{N,j}=\binom Nj\langle v,T^j(I-T)^{N-j}v\rangle\ge0,
 \qquad \sum_{j=0}^Nw_{N,j}=1.
\]
\end{proposition}
\begin{proof}
Each weight is nonnegative because its defining polynomial
$\binom Nj q^j(1-q)^{N-j}$ is nonnegative on $[0,1]$.
The binomial identity for the commuting pair $(T,I-T)$
shows that the sum of the operators is $I$, and normalization
then gives total mass one.
\end{proof}

The first two difference orders already have recognizable content:
$u_n\ge u_{n+1}$ and $u_n-2u_{n+1}+u_{n+2}\ge0$.
Thus positive models produce decreasing, discretely convex sequences.
Higher orders impose further conditions; neither monotonicity nor
convexity alone characterizes a Hausdorff moment sequence.

The moments also give positive Hankel forms. For arbitrary complex
scalars $a_0,\ldots,a_d$,
\[
 \sum_{i,j=0}^{d}\overline{a_i}a_j u_{i+j}
 =\left\|\sum_{j=0}^{d}a_jT^jv\right\|^2\ge0.
\]
Here self-adjointness is used to identify
$\langle T^iv,T^jv\rangle$ with $u_{i+j}$.
The identities explain how vector positivity is visible through
finite scalar tests. Their converse requires a representation
argument: checking finitely many particular rows does not establish
all orders, and a formal interface with positive fields does not
prove that those fields can be supplied for the intended sequence.

\PaperPage{11}{1.5 The concrete multiplication model}
\label{hilb:multiplication}
Let $\sigma$ be a Borel probability measure on $I=[0,1]$. Consider
the complex Hilbert space
\[
 \mathcal H=L^2(I,\sigma;\mathbb C),\qquad
 \langle f,g\rangle=\int_I\overline{f(q)}g(q)\,d\sigma(q).
\]
Functions here are equivalence classes modulo equality almost
everywhere. Multiplication by the bounded coordinate function
preserves that equivalence and defines a bounded linear operator:
\[
 (Tf)(q)=qf(q),\qquad v(q)=1.
\]
Since $|q|\le1$, one has $\|Tf\|_2\le\|f\|_2$.
The constant function belongs to the space and has norm one.

\begin{proposition}\label{hilb:measure-realization}
The triple $(L^2(\sigma),T,1)$ is a positive Hilbert-vector model.
For every $n\ge0$ it realizes the moments of $\sigma$:
\[
 \langle 1,T^n1\rangle=\int_Iq^n\,d\sigma(q).
\]
\end{proposition}
\begin{proof}
The coordinate $q$ is real, so
$\langle Tf,g\rangle=\langle f,Tg\rangle$ for all $f,g$.
A bounded operator defined on the whole space with this symmetry
is self-adjoint. Moreover,
$\langle f,Tf\rangle=\int q|f|^2\,d\sigma$ lies between zero
and $\|f\|_2^2$. Induction gives $T^n1=q^n$, proving the
moment formula.
\end{proof}

If $z\notin[1,\infty)$, the continuous function
$q\mapsto1-zq$ has no zero on the compact interval $I$.
Its reciprocal is therefore bounded. Multiplication by that
reciprocal is a bounded two-sided inverse of $I-zT$, giving
\[
 \langle1,(I-zT)^{-1}1\rangle
 =\int_I\frac{d\sigma(q)}{1-zq}.
\]
Thus the vector resolvent is exactly an integral transform, including
at $z=0$. No choice of basis or eigenvector expansion is involved,
and no density assumption on $\sigma$ is needed. Atomic, singular,
and absolutely continuous measures are treated identically.

On a compact subset of this slit domain, $|1-zq|$ has a
positive lower bound uniform in both $z$ and $q$. Difference
quotients and their limits are consequently dominated by an
integrable constant. Differentiation under the integral is valid,
proving holomorphy directly throughout that domain.

This model is the concrete Hilbert-space interpretation of the
probability measure constructed later from positive Bernstein rows.
The audited project proves that measure construction and its scalar
integral consequences. The present operator packaging is mathematical
exposition; it is not an additional Lean construction already covered
by the source audit.

\PaperPage{12}{1.6 Polynomial density and uniqueness of the measure}
\label{hilb:density}
Two standard approximation principles make the multiplication model
particularly precise. First, complex polynomials restricted to $I$
are uniformly dense in $C(I,\mathbb C)$ by Stone--Weierstrass.
Second, continuous functions are dense in $L^2(I,\sigma)$ for a
finite Borel measure on this compact metric space \cite{rudin}.
The latter follows from regularity: approximate indicators by
continuous functions between a compact subset and an open
superset, and then approximate simple functions.

Since $\sigma(I)=1$, uniform approximation implies approximation
in the Hilbert norm:
\[
 \|f-p\|_{L^2(\sigma)}\le\|f-p\|_\infty
       \qquad(f\in C(I),\ p\text{ a polynomial}).
\]
It follows that the polynomial classes are dense in $L^2(\sigma)$.
But $p(T)1=p(q)$, so
\[
 \overline{\operatorname{span}}\{1,T1,T^21,\ldots\}
     =L^2(I,\sigma).
\]
The distinguished vector in the multiplication model is cyclic,
even when the measure has atoms or is concentrated on a proper
closed subset of the interval.

\begin{proposition}\label{hilb:moment-uniqueness}
A Borel probability measure on $[0,1]$ is uniquely determined by
all its monomial moments.
\end{proposition}
\begin{proof}
Suppose $\sigma$ and $\tau$ have the same moments. Linearity gives
equal integrals for every polynomial. If $p_k\to f$ uniformly, then
\[
 \left|\int f\,d\sigma-\int p_k\,d\sigma\right|
 \le\|f-p_k\|_\infty,
\]
and the same bound holds for $\tau$. The two integrals therefore
agree for every continuous $f$. Uniqueness in the Riesz
representation theorem implies $\sigma=\tau$.
\end{proof}

All moments are essential in this statement. A finite collection
determines integrals only on a finite-dimensional polynomial
space and generally permits many representing measures.
For instance, the unit mass at $1/2$ and normalized Lebesgue
measure on $[0,1]$ have the same zeroth and first moments.
Their second moments are $1/4$ and $1/3$, respectively. Even
exact agreement at those first two orders leaves different
resolvents at the next Taylor coefficient.
The eventual construction of the actual measure thus has two
distinct features: existence follows from all-order positivity,
while exact moment identification fixes its scalar resolvent
uniquely. Neither fact supplies a terminating decision procedure
for the all-order positivity condition.

\PaperPage{13}{1.7 Spectral measures and the visible operator}
\label{hilb:visible}
For a bounded self-adjoint positive contraction $T$, the spectral
theorem provides a projection-valued measure $E$ supported in
$[0,1]$. A unit vector $v$ turns it into a scalar probability measure:
\[
 \sigma_v(B)=\langle v,E(B)v\rangle=\|E(B)v\|^2.
\]
The spectral calculus then gives, for continuous $f$,
\[
 \langle v,f(T)v\rangle=\int_{[0,1]} f(q)\,d\sigma_v(q).
\]
In particular, the moments and resolvent are those of $\sigma_v$.
This recovers the integral description from an abstract positive
model; the preceding multiplication construction provides the
converse direction.

The cyclic restriction is determined up to unitary equivalence.
Indeed, define on polynomials
\[
 U\colon p\longmapsto p(T)v.
\]
The moment identities imply
$\|p(T)v\|^2=\int|p(q)|^2\,d\sigma_v(q)$.
Thus $U$ respects the null classes in $L^2(\sigma_v)$ and is
isometric on the polynomial subspace. Density extends it to a
unitary map from $L^2(\sigma_v)$ onto $\mathcal H_v$.
The identity $U(qp)=T\,U(p)$ on polynomials extends by
boundedness, so $U$ intertwines multiplication by $q$ with
$T|_{\mathcal H_v}$.

There is no corresponding determination of the unobserved
orthogonal complement. For any other positive contraction $A$,
the enlarged model
\[
 (\mathcal H\oplus\mathcal K,\ T\oplus A,\ v\oplus0)
\]
has exactly the same scalar moments and resolvent as the original
one. The spectrum of $A$ may add points not present in the
cyclic restriction. Consequently a scalar pole or support
calculation must be interpreted relative to the represented
vector, not automatically as a statement about an arbitrary
ambient operator.

This distinction matters when using operator language in a zero
problem. Constructing some self-adjoint operator is not enough:
one must specify the distinguished vector, prove the exact
moment identities, and identify the associated analytic function.
The actual-function calibration in the RH application performs
that final identification at the scalar level. It is not
replaced by the existence of the abstract spectral representation.

\PaperPage{14}{1.8 Elementary models and limits of the assumptions}
\label{hilb:examples}
Explicit examples distinguish contractivity, positivity, and
actual-function identification without introducing analytic
complications.

\begin{example}[A single spectral atom]
On $\mathbb C$, let $T=qI$, $v=1$, and $0\le q\le1$. Then
\[
 u_n=q^n,\qquad F(z)=\frac1{1-qz},\qquad
 w_{N,j}=\binom Nj q^j(1-q)^{N-j}.
\]
The corresponding measure is $\delta_q$. At $q=0$ the
resolvent is constant; at $q=1$ it has a pole at $z=1$.
Thus the guaranteed unit disk need not be the maximal domain.
\end{example}

\begin{example}[Two atoms]
Let $T=\operatorname{diag}(q_1,q_2)$ on $\mathbb C^2$ and
$v=(\sqrt t,\sqrt{1-t})$, where $q_1,q_2,t\in[0,1]$.
Then
\[
 u_n=tq_1^n+(1-t)q_2^n,\qquad
 F(z)=\frac{t}{1-q_1z}+\frac{1-t}{1-q_2z}.
\]
The scalar measure is $t\delta_{q_1}+(1-t)\delta_{q_2}$.
If $t=1$, the second coordinate is invisible to $F$, even
when $q_2\ne q_1$.
For distinct $q_1,q_2$ and $0<t<1$, however, $v$ and $Tv$
are linearly independent: their coordinate determinant is
$\sqrt{t(1-t)}(q_2-q_1)\ne0$. Hence $v$ is cyclic in this
case. Both atoms are visible, illustrating directly why positive
weights and distinct support points enter the cyclic reconstruction.
\end{example}

A self-adjoint contraction need not be positive. For $T=-I$
and $v=1$, the coefficients are $c_n=(-1)^n$ and
$F(z)=(1+z)^{-1}$. The coefficient bound and disk analyticity
hold, but $c_1<0$, so these are not moments of a positive
measure on $[0,1]$. The positivity assumption cannot be
recovered from the uniform coefficient estimate.

Self-adjointness also matters for the reconstruction statement.
Let $S$ be the unilateral shift on $\ell^2(\mathbb N)$ and
$v=e_0$. The orbit is orthonormal, giving
\[
 \langle e_0,S^ne_0\rangle=
 \begin{cases}1,&n=0,\\0,&n>0.\end{cases}
 \qquad F(z)=1\quad(|z|<1).
\]
The vector is cyclic, yet this is the same scalar generating
function as the zero operator on $\mathbb C$. The two
operators are not unitarily equivalent. Scalar orbit moments
of a general contraction do not determine mixed orbit inner
products. The positive self-adjoint setting supplies exactly
that extra structure; the actual arithmetic application still
requires its separate calibration even after positivity is known.
